
\documentclass[12pt,a4paper]{article}

\usepackage{xepersian}
\settextfont{XB Niloofar}
\setlatintextfont{Times New Roman}
\usepackage{amsmath,amssymb}
\usepackage{graphicx}
\usepackage{geometry}
\geometry{margin=2.5cm}
\usepackage{setspace}
\onehalfspacing
\usepackage{hyperref}

\title{%
گزارش پژوهشی\\
\large مدل‌سازی و بهبود پیش‌بینی خوانش چهار کیوبیت با استفاده از روش‌های یادگیری ماشین و کرنل‌های الهام‌گرفته از کوانتوم
}
\author{ }
\date{}

\begin{document}
\maketitle

\section*{چکیده}
در این پروژه، مسئله‌ی پیش‌بینی حالت کامل یک سامانه‌ی چهار کیوبیتی بر اساس داده‌های خوانش کوادریچر مورد بررسی قرار گرفته است. داده‌ها شامل اندازه‌گیری‌های مختلط برای هر کیوبیت بوده و هدف، نگاشت هر نمونه‌ی اندازه‌گیری به یکی از ۱۶ حالت ممکن سامانه‌ی چهار کیوبیتی است. با وجود استفاده از روش‌های مختلف یادگیری ماشین کلاسیک، مهندسی ویژگی و کرنل‌های الهام‌گرفته از نگاشت‌های کوانتومی نظیر \lr{ZZ feature map}، دقت پیش‌بینی چندکلاسه در بهترین حالت به حدود ۳۸ تا ۳۹ درصد محدود شده است. در این گزارش، مسئله از دید فیزیکی و آماری تشریح شده، روش‌های پردازش داده و مدل‌سازی بررسی می‌شوند و در نهایت مسیرهای فیزیکی، آزمایشگاهی و الگوریتمی برای بهبود دقت مدل ارائه می‌گردد.

\section{مقدمه و شرح مسئله}
خوانش حالت کیوبیت‌ها یکی از چالش‌های اساسی در سامانه‌های محاسبات کوانتومی مبتنی بر پلتفرم‌هایی نظیر کیوبیت‌های ابررسانا یا یون‌های به‌دام‌افتاده است. در این سامانه‌ها، اطلاعات حالت کوانتومی به‌طور مستقیم قابل اندازه‌گیری نیست و از طریق فرآیندهای خوانش غیرمستقیم استخراج می‌شود. این فرآیندها معمولاً منجر به سیگنال‌هایی در فضای کوادریچر می‌شوند که به‌صورت نقاطی در صفحه‌ی مختلط \lr{(I,Q)} نمایش داده می‌شوند.

در این تحقیق، داده‌های مربوط به چهار کیوبیت مورد بررسی قرار گرفته‌اند. هر نمونه‌ی داده شامل چهار عدد مختلط است که هر یک متناظر با خوانش یک کیوبیت می‌باشد. برچسب هر نمونه یکی از ۱۶ حالت ممکن \lr{|q_4 q_3 q_2 q_1\rangle} است. هدف اصلی، طراحی یک مدل طبقه‌بندی است که بتواند با دریافت این چهار خوانش مختلط، حالت کامل سامانه‌ی چهار کیوبیتی را پیش‌بینی کند.

\section{ساختار داده‌ها و تحلیل آماری اولیه}
هر مجموعه‌داده از ۱۶ فایل تشکیل شده است که هر فایل متناظر با یک حالت چهارکیوبیتی خاص می‌باشد. در هر فایل، برای هر کیوبیت ۱۰۰۰ نمونه‌ی اندازه‌گیری ثبت شده است. با ترکیب این داده‌ها، مجموعه‌داده‌ای شامل ۱۶۰۰۰ نمونه ایجاد شده است.

تحلیل‌های بصری اولیه شامل ترسیم نمودارهای پراکندگی در صفحه‌ی \lr{IQ}، هیستوگرام‌ها و توابع توزیع تجمعی انجام شد. این تحلیل‌ها نشان داد که توزیع داده‌های مربوط به حالات مختلف به‌شدت هم‌پوشان هستند. حتی با تقریب توزیع‌ها به شکل گاوسی و نمایش بیضی‌های کوواریانس، جدایش خطی یا حتی غیرخطی ساده بین کلاس‌ها امکان‌پذیر نبود. این رفتار با گزارش‌های تجربی رایج در خوانش کیوبیت‌ها هم‌خوانی دارد.

\section{مهندسی ویژگی و پیش‌پردازش داده}
در مرحله‌ی مهندسی ویژگی، دو نمایش اصلی برای داده‌های مختلط در نظر گرفته شد. در نمایش نخست، هر عدد مختلط به دو مولفه‌ی حقیقی و موهومی تجزیه شد که منجر به بردار ویژگی ۸بعدی می‌گردد. در نمایش دوم که به‌عنوان \lr{Angle Encoding} شناخته می‌شود، هر عدد مختلط به دامنه و فاز متناظر آن نگاشت شد. این نمایش از نظر فیزیکی معنادارتر بوده و به‌ویژه در ترکیب با کرنل‌های مبتنی بر توابع تناوبی مزیت دارد.

به‌منظور کاهش اثر داده‌های پرت و ناپایداری آماری، روش‌های مختلف نرمال‌سازی بررسی شدند. نتایج نشان داد که استفاده از \lr{RobustScaler} نسبت به \lr{StandardScaler} برای داده‌های نویزی آزمایشگاهی مناسب‌تر است.

\section{مدل‌سازی و الگوریتم‌های یادگیری ماشین}
مدل پایه‌ی مورد استفاده، ماشین بردار پشتیبان \lr{(SVM)} بود. در ابتدا، کرنل‌های استاندارد نظیر \lr{RBF} برای طبقه‌بندی تک‌کیوبیتی و چندکلاسه مورد استفاده قرار گرفتند. در حالت تک‌کیوبیتی، دقت‌های نسبتاً بالایی برای برخی کیوبیت‌ها به‌دست آمد، اما در حالت چندکلاسه، دقت به‌طور قابل توجهی کاهش یافت.

به‌منظور لحاظ کردن هم‌بستگی‌های بین ویژگی‌ها، یک کرنل سفارشی الهام‌گرفته از \lr{ZZ feature map} معرفی شد. این کرنل شامل ترم‌های تک‌ویژگی و ترم‌های جفتی بود که به‌صورت توابع کسینوسی تعریف می‌شدند. تنظیم پارامترهای این کرنل از طریق جستجوی شبکه‌ای انجام شد.

\section{نتایج و محدودیت‌ها}
بهترین دقت حاصل برای مسئله‌ی ۱۶کلاسه حدود ۳۸ تا ۳۹ درصد بود. این نتیجه نشان می‌دهد که اگرچه کرنل‌های الهام‌گرفته از کوانتوم می‌توانند اطلاعات بیشتری از ساختار داده‌ها استخراج کنند، اما هم‌پوشانی ذاتی توزیع‌ها و نویز اندازه‌گیری محدودیت اصلی باقی می‌ماند.

\section{مسیرهای پیشنهادی برای بهبود}
از منظر فیزیکی و آزمایشگاهی، بهبود تفکیک‌پذیری خوانش، کاهش \lr{cross-talk} و افزایش نسبت سیگنال به نویز می‌تواند نقش کلیدی ایفا کند. از منظر مدل‌سازی، استفاده از مدل‌های مولد، روش‌های \lr{ensemble} و مدل‌های گرافی پیشنهاد می‌شود. در نهایت، بهره‌گیری از الگوریتم‌های کوانتومی واقعی نظیر \lr{QSVM} یا طبقه‌بندهای واریاسیونی کوانتومی می‌تواند مسیر آینده‌ی این تحقیق باشد.


\section{تحلیل هم‌پوشانی ذاتی و وابستگی خوانش چندکیوبیتی}
به‌منظور تشخیص این‌که محدودیت دقت مدل ناشی از هم‌پوشانی ذاتی توزیع‌های خوانش است یا وابستگی و \lr{cross-talk} بین کیوبیت‌ها، یک تحلیل داده‌محور با الهام از معیارهای معرفی‌شده در مقاله‌ی
\lr{Multiplexed qubit readout quality metric beyond assignment fidelity}
انجام شد. تأکید این تحلیل بر استفاده‌ی صرف از داده‌های موجود خوانش بوده و هیچ‌گونه دسترسی به پارامترهای فیزیکی آزمایش فرض نشده است.

\subsection{جدایش‌پذیری ذاتی تک‌کیوبیتی}
در گام نخست، برای هر کیوبیت یک مدل طبقه‌بندی دودویی تنها بر اساس خوانش همان کیوبیت آموزش داده شد. دقت این مدل‌ها را می‌توان به‌عنوان سقف اطلاعاتی جدایش‌پذیری ذاتی هر کیوبیت در نظر گرفت. نتایج به‌صورت زیر به‌دست آمد:
\begin{center}
\begin{tabular}{c c}
\hline
کیوبیت & دقت جدایش ذاتی \\
\hline
کیوبیت ۱ & 0.713 \\
کیوبیت ۲ & 0.818 \\
کیوبیت ۳ & 0.802 \\
کیوبیت ۴ & 0.814 \\
\hline
\end{tabular}
\end{center}

این نتایج نشان می‌دهند که حتی در بهترین حالت و با استفاده از اطلاعات خوانش همان کیوبیت، خطای ذاتی قابل توجهی وجود دارد. اگر فرض شود خطاهای خوانش کیوبیت‌ها مستقل هستند، سقف دقت پیش‌بینی حالت چهارکیوبیتی را می‌توان به‌صورت تقریبی برابر حاصل‌ضرب دقت‌های تک‌کیوبیتی برآورد کرد که عددی در حدود ۳۸ درصد به‌دست می‌دهد. این مقدار با بهترین دقت مشاهده‌شده در مدل چندکلاسه هم‌خوانی دارد.

\subsection{تحلیل وابستگی بین‌کیوبیتی (Cross-talk)}
در گام دوم، وابستگی خوانش کیوبیت $i$ به حالت کیوبیت $j$ بررسی شد. برای این منظور، توزیع خوانش کیوبیت $i$ مشروط به $q_j=0$ و $q_j=1$ مقایسه گردید. معیارهایی نظیر جابه‌جایی میانگین و فاصله‌ی ماهالانوبیس بین این دو توزیع محاسبه شدند.

نتایج نشان دادند که عناصر روی قطر ماتریس وابستگی (یعنی $i=j$) به‌طور قابل توجهی بزرگ‌تر از عناصر خارج از قطر هستند، در حالی که وابستگی‌های بین‌کیوبیتی مقدار بسیار کوچکی دارند. اگرچه برخی از این وابستگی‌های ضعیف از نظر آماری معنی‌دار هستند، اما اندازه‌ی اثر آن‌ها از نظر فیزیکی ناچیز است. این موضوع نشان می‌دهد که \lr{cross-talk} بین خوانش‌ها نقش غالبی در محدودیت دقت مدل ایفا نمی‌کند.

\subsection{نتیجه‌گیری تحلیلی}
این تحلیل نشان می‌دهد که عامل اصلی محدودکننده‌ی دقت پیش‌بینی حالت چهارکیوبیتی، هم‌پوشانی ذاتی توزیع‌های خوانش تک‌کیوبیتی است، نه وابستگی شدید بین خوانش کیوبیت‌ها. بنابراین، افزایش قابل توجه دقت مدل تنها از طریق پیچیده‌تر کردن الگوریتم‌های یادگیری ماشین ممکن نیست و نیازمند بهبودهای فیزیکی در فرآیند خوانش و افزایش تفکیک‌پذیری اندازه‌گیری‌ها است.


\section{جمع‌بندی}
در این گزارش، مسئله‌ی پیش‌بینی حالت کامل چهار کیوبیت با استفاده از داده‌های خوانش کوادریچر بررسی شد. نتایج نشان داد که محدودیت دقت مدل بیش از آنکه ناشی از ضعف الگوریتم‌های یادگیری ماشین باشد، ریشه در ماهیت فیزیکی فرآیند اندازه‌گیری دارد. مسیرهای پیشنهادی ارائه‌شده می‌توانند مبنایی برای ادامه‌ی تحقیق در حوزه‌های آزمایشگاهی و الگوریتمی باشند.

\end{document}
